% Twin claim-proof file for row claim_3_22 -- "SigmaSeparationSymmetric"
% under the `collider_side_aware` refactor (root: def_3_15).
%
% This is the REFACTOR TWIN of `claim_3_22_proof_SigmaSeparationSymmetric.tex`.
% The original sibling file stays in place untouched until Phase 7 cleanup,
% at which point the cleanup script renames THIS file over the original.
% The two files render the same mathematical claim with the same overall
% proof strategy; the side-aware refactor changes only the *per-step
% arrowhead-contribution reading* used in beats (III.a) and (III.b) below.
% The LN block being proven is unchanged.
%
% Source of the proof framework: the original tex proof
% `tex/claim_3_22_proof_SigmaSeparationSymmetric.tex` (beats I, II, III.c,
% IV, V carried over essentially verbatim; beats III.a and III.b rewritten
% to thread the side-aware reading through explicitly).
%
% Authoritative references for the side-aware reading:
%   - `tex/def_3_15_CollidersAndNon.tex`, paragraphs "Encoding note: side-
%     aware reading via a typed walk-step representation" and "Treatment
%     of directed self-loops"; addition tag
%     `[collider_side_aware_walkstep_predicates]`.
%   - `leanification/Chapter3_GraphTheory/Section3_1/CDMG.lean:376`
%     (`hL_irrefl : ∀ ⦃s : Sym2 Node⦄, s ∈ L → ¬ s.IsDiag`) — closes the
%     L-disjunct vacuously at a directed self-loop.
%   - `leanification/deviations.json`, manager-accepted deviation
%     `collider_side_aware_at_self_loops` — fixes the canonical
%     disambiguation at directed-self-loop steps.

\documentclass[main]{subfiles}

\begin{document}
% ref: claim_3_22 (proof, with statement restated for readability; refactor twin)
% title: SigmaSeparationSymmetric
\def\rowref{claim\_3\_22}
\phantomsection\label{claim_3_22}

% Local macros for the side-aware encoding helpers used in beats (III.a)
% and (III.b). Scoped to this file via `\providecommand` (no-op if the
% identifier is later upstreamed into the shared preamble). Match the
% Lean identifiers from `def_3_15` / `def_3_16` for grep-readability.
\providecommand{\HeadAtTarget}{\mathrm{HeadAtTarget}}
\providecommand{\HeadAtSource}{\mathrm{HeadAtSource}}
\providecommand{\IsCollider}{\mathrm{IsCollider}}
\providecommand{\IsNonCollider}{\mathrm{IsNonCollider}}
\providecommand{\HasBlockingLeftSlot}{\mathrm{HasBlockingLeftSlot}}
\providecommand{\HasBlockingRightSlot}{\mathrm{HasBlockingRightSlot}}
\providecommand{\WalkStep}{\texttt{WalkStep}}
\providecommand{\Node}{\texttt{Node}}
\providecommand{\IsDiag}{\mathrm{IsDiag}}
\providecommand{\Sym}{\mathrm{Sym}}

% --- statement (restated, synced verbatim from
% --- claim_3_22_statement_SigmaSeparationSymmetric.tex) -------------------
\begin{Lem}[$\sigma$-separation is symmetric]\label{lem:sigma_separation_symmetric}
    Let $G = (J, V, E, L)$ be a CDMG (def \ref{def-cdmg}); throughout we use the notation of def \ref{not-cdmg}, the walk definition of def \ref{def:walks}, item~i., the $\sigma$-blocking predicate (``$\pi$ is $C$-$\sigma$-blocked'') of def \ref{def:sigma_blocking}, and the $i\sigma$-/$\sigma$-separation predicates of def \ref{def:sigma_separation}. Suppose
    \[
        J = \emptyset,
    \]
    so that, by def \ref{def:sigma_separation}, item~4, the $\sigma$-separation abbreviation
    \[
        A \sPerp_G B \given C \;\; := \;\; A \isPerp_G B \given C
    \]
    is in scope (and, as recorded there, $A \sPerp_G B \given C$ unfolds under $J = \emptyset$ to: for every integer $n \ge 0$ and every walk $\pi = (v_0, a_0, v_1, \dots, a_{n-1}, v_n)$ in $G$ in the sense of def \ref{def:walks}, item~i., with $v_0 \in A$ and $v_n \in B$, the walk $\pi$ is $C$-$\sigma$-blocked in the sense of def \ref{def:sigma_blocking}).
    Let $A, B, C \ins V$ be subsets of nodes --- equivalently $A, B, C \ins J \cup V$, since $J = \emptyset$; $A$, $B$, $C$ are \emph{not} assumed to be pairwise disjoint, nor non-empty.

    Then $\sigma$-separation in $G$ is symmetric in its first two arguments:
    \[
        A \sPerp_G B \given C \iff B \sPerp_G A \given C.
    \]
\end{Lem}

% --- proof ------------------------------------------------------------------
\begin{proof}
The proof has three main beats:
\begin{enumerate}
    \item[(I)] unfold $\sPerp_G$ via def \ref{def:sigma_separation}, item~4, into the universal-over-walks statement of def \ref{def:sigma_separation}, item~1, and reduce the biconditional to its forward implication (the converse follows by swapping $A$ and $B$ in the same argument);
    \item[(II)] construct an explicit walk-reversal involution $\pi \mapsto \pi^*$ that is a bijection between walks $v_0 \to v_n$ and walks $v_n \to v_0$ in $G$ in the sense of def \ref{def:walks}, item~i.; and
    \item[(III)] show that the $C$-$\sigma$-blocking predicate of def \ref{def:sigma_blocking} is invariant under that reversal, so that the universal-over-walks statement is symmetric in $A$ and $B$.
\end{enumerate}
The substantive content lives in beats (II) and (III): once walk reversal is set up as an involution that maps $C$-$\sigma$-blocked walks to $C$-$\sigma$-blocked walks (and back), the biconditional drops out by direct chase.

The mathematical content of this proof is identical to that of the sibling original tex proof file
\texttt{tex/claim\_3\_22\_proof\_SigmaSeparationSymmetric.tex}; the differences with the original
lie entirely in beats (III.a) and (III.b), where the per-step arrowhead-contribution reading is the
\emph{side-aware} one of \refrow{def_3_15}'s addition tag
\texttt{[collider\_side\_aware\_walkstep\_predicates]} -- reading per-step head contributions off
the typed walk-step's constructor-tagged channel and walk-traversal direction (forward-$E$ /
backward-$E$ / bidirected-$L$), with an opposite-side $L$-channel re-firing disjunct that catches
writing-mirror pairs and that is vacuously false at directed self-loops by \refrow{def_3_1}'s
\texttt{hL\_irrefl} restriction.  The side-aware reading agrees pointwise with the literal
stored-pair ``edge into $v_k$'' test of def \ref{def-edge-relations}, item~ii.\ on all walks that
traverse no directed self-loop step, and refines it at directed-self-loop steps (cf.\ \refrow{def_3_15}'s
``Treatment of directed self-loops'' paragraph; manager-accepted deviation
\texttt{collider\_side\_aware\_at\_self\_loops}).  The walk-reversal involution and the $\sigma$-blocking
disjunctive structure are unchanged, so beats (I), (II), (III.c), (IV), and (V) carry over from the
original modulo a clarifying re-statement of beat (II)'s case-clauses in terms of the typed walk-
step constructors.

\medskip\noindent\textbf{(I) Unfold and reduce to one direction.}
Throughout, fix $G = (J, V, E, L)$ a CDMG (def \ref{def-cdmg}) with $J = \emptyset$, and $A, B, C \ins V = J \cup V$. By def \ref{def:sigma_separation}, item~4, the abbreviation $A \sPerp_G B \given C$ is in scope and unfolds (under $J = \emptyset$, so $J \cup B = B$ and $J \cup A = A$) to
\[
    A \sPerp_G B \given C
    \quad\iff\quad
    \forall\, n \ge 0, \;\forall\, \text{walks } \pi = (v_0, a_0, \dots, a_{n-1}, v_n) \text{ in } G
    \text{ with } v_0 \in A,\, v_n \in B :\;
    \pi \text{ is } C\text{-}\sigma\text{-blocked},
\]
where ``walk in $G$'' is in the sense of def \ref{def:walks}, item~i., and ``$C$-$\sigma$-blocked'' is in the sense of def \ref{def:sigma_blocking}. Symmetrically $B \sPerp_G A \given C$ unfolds to the analogous universal over walks $\pi'$ in $G$ with $v_0' \in B$ and $v_n' \in A$.

The biconditional we have to prove is symmetric in the roles of $A$ and $B$: the right-hand side $B \sPerp_G A \given C$ is obtained from the left-hand side $A \sPerp_G B \given C$ by literally swapping the symbols $A$ and $B$. It therefore suffices to prove the forward implication
\[
    A \sPerp_G B \given C \;\implies\; B \sPerp_G A \given C; \tag{$\star$}
\]
the converse implication then follows by applying $(\star)$ to the pair $(B, A)$ in place of $(A, B)$. From now on we assume $A \sPerp_G B \given C$ and aim to show $B \sPerp_G A \given C$.

\medskip\noindent\textbf{(II) Walk reversal.}
We construct, for every $n \ge 0$ and every walk $\pi = (v_0, a_0, v_1, a_1, \dots, a_{n-1}, v_n)$ in $G$ (def \ref{def:walks}, item~i.), a \emph{reversed walk} $\pi^*$ from $v_n$ to $v_0$ in $G$. The construction is per slot, by the walk constraint of def \ref{def:walks}, item~i.

For $i \in \{0, 1, \dots, n - 1\}$, the walk constraint at slot $i$ of $\pi$ states
\[
    \bigl(a_i = (v_i, v_{i+1}) \;\land\; a_i \in E \cup L\bigr)
    \;\lor\;
    \bigl(a_i = (v_{i+1}, v_i) \;\land\; a_i \in E\bigr).
\]
The slot $i$ of $\pi$ contributes an edge $a_i \in E \cup L$ to the walk; the typed encoding records, in addition to the stored ordered pair, the \emph{channel and walk-traversal direction} via one of three constructor tags --- one tag per disjunct of the walk constraint, plus a third tag for the bidirected $L$-channel:
\begin{itemize}
    \item[(F$_E$)] The first disjunct's $E$-instance ($a_i = (v_i, v_{i+1}) \in E$) corresponds to the constructor tag \texttt{forwardE}, encoding a directed graph-edge of $G$ with tail $v_i$ and head $v_{i+1}$, traversed forwards $v_i \to v_{i+1}$ on $\pi$.
    \item[(B$_E$)] The second disjunct ($a_i = (v_{i+1}, v_i) \in E$) corresponds to the constructor tag \texttt{backwardE}, encoding a directed graph-edge with tail $v_{i+1}$ and head $v_i$, traversed backwards $v_i \leftarrow v_{i+1}$ on $\pi$ (i.e.\ in the walk-traversal direction $v_i \to v_{i+1}$).
    \item[(L)] The first disjunct's $L$-instance ($a_i = (v_i, v_{i+1}) \in L$, equivalently $s(v_i, v_{i+1}) \in L$ with $L$ stored as $\Sym_2$-typed unordered pairs per def \ref{def-cdmg}) corresponds to the constructor tag \texttt{bidir}, encoding a bidirected graph-edge with arrowheads at both endpoints.
\end{itemize}
We define the reversed edge $a_i^*$ at slot $n - 1 - i$ of $\pi^*$ as follows, by case-analysis on which constructor tag is recorded at slot $i$ of $\pi$:
\begin{itemize}
    \item[(F$_E$)] If the slot is a \texttt{forwardE} step with $a_i = (v_i, v_{i+1}) \in E$, set
        \[
            a_i^* := (v_i, v_{i+1}) \in E,
        \]
        the \emph{same} directed graph-edge, but now occupying slot $n - 1 - i$ of $\pi^*$ between the nodes $v_{i+1}$ (left) and $v_i$ (right). The walk constraint of def \ref{def:walks}, item~i.\ at slot $n - 1 - i$ of $\pi^*$ is then satisfied via its \emph{second} disjunct ($a_i^* = (\text{right-node}, \text{left-node}) = (v_i, v_{i+1}) \in E$), i.e.\ the constructor tag at slot $n - 1 - i$ of $\pi^*$ is \texttt{backwardE}: a directed graph-edge that was traversed forwards on $\pi$ is read backwards on $\pi^*$.
    \item[(B$_E$)] If the slot is a \texttt{backwardE} step with $a_i = (v_{i+1}, v_i) \in E$, set
        \[
            a_i^* := (v_{i+1}, v_i) \in E,
        \]
        the same directed graph-edge, now occupying slot $n - 1 - i$ of $\pi^*$ between $v_{i+1}$ (left) and $v_i$ (right). The walk constraint at slot $n - 1 - i$ of $\pi^*$ is satisfied via its \emph{first} disjunct ($a_i^* = (\text{left-node}, \text{right-node}) = (v_{i+1}, v_i) \in E \subseteq E \cup L$), so the constructor tag at slot $n - 1 - i$ of $\pi^*$ is \texttt{forwardE}: a directed graph-edge traversed backwards on $\pi$ becomes one traversed forwards on $\pi^*$.
    \item[(L)] If the slot is a \texttt{bidir} step with $a_i \in L$ stored as $s(v_i, v_{i+1}) \in L$, set
        \[
            a_i^* := s(v_{i+1}, v_i) \in L = s(v_i, v_{i+1}) \in L,
        \]
        using the canonical $\Sym_2$ identification $s(v_{i+1}, v_i) = s(v_i, v_{i+1})$ (def \ref{def-cdmg}: $L$ is stored as unordered pairs over $V$). The walk constraint at slot $n - 1 - i$ of $\pi^*$ is satisfied via its first disjunct ($a_i^* \in L \subseteq E \cup L$, with left-node $v_{i+1}$ and right-node $v_i$), so the constructor tag at slot $n - 1 - i$ of $\pi^*$ is again \texttt{bidir}: a bidirected graph-edge is bidirected in both directions.
\end{itemize}
At the constructor-tag level the reversal acts as the involution \texttt{forwardE} $\leftrightarrow$ \texttt{backwardE}, \texttt{bidir} $\leftrightarrow$ \texttt{bidir}. The reversed walk is then the tuple
\[
    \pi^* \;:=\; (v_n, \; a_{n-1}^*, \; v_{n-1}, \; a_{n-2}^*, \; \dots, \; a_1^*, \; v_1, \; a_0^*, \; v_0),
\]
i.e.\ the same node sequence read in reverse, paired with the slot-by-slot edge construction above. By the per-slot verification, $\pi^*$ satisfies the walk constraint of def \ref{def:walks}, item~i.\ at each of its $n$ slots, so $\pi^*$ is a walk in $G$ from $v_n$ to $v_0$.

\smallskip\noindent\emph{Involution.}
Each of the three cases above is its own inverse at the constructor-tag level: applying the construction once flips the constructor tag (\texttt{forwardE} $\leftrightarrow$ \texttt{backwardE}; \texttt{bidir} $\leftrightarrow$ \texttt{bidir}) and rewrites the stored pair / $\Sym_2$-witness accordingly, and applying it a second time flips it back. Concretely, in case (F$_E$) the reverse of $\pi^*$ at slot $n - 1 - (n - 1 - i) = i$ is again $a_i = (v_i, v_{i+1}) \in E$ (we read $\pi^*$'s slot $n - 1 - i$, whose stored edge is $(v_i, v_{i+1}) \in E$ with constructor tag \texttt{backwardE}, in case (B$_E$) of the reversal construction applied to $\pi^*$, and recover the first-disjunct/\texttt{forwardE} configuration of the original $\pi$). Cases (B$_E$) and (L) are analogous (the L case using the $\Sym_2$ identification $s(v_{i+1}, v_i) = s(v_i, v_{i+1})$ canonically). Hence $(\pi^*)^* = \pi$ on the nose, so $\pi \mapsto \pi^*$ is an involution on the set of walks in $G$.

\smallskip\noindent\emph{Bijection between walk-spaces.}
For fixed source-set $A$ and target-set $B$, the involution restricts to a bijection between
\[
    \mathcal{W}_{A \to B} \;:=\; \{\,\pi \text{ walk in } G : v_0 \in A,\, v_n \in B\,\}
    \quad\text{and}\quad
    \mathcal{W}_{B \to A} \;:=\; \{\,\pi^* \text{ walk in } G : v_0^* \in B,\, v_n^* \in A\,\},
\]
since $\pi \in \mathcal{W}_{A \to B}$ has $v_0 \in A,\, v_n \in B$ and $\pi^*$ has its start-node equal to $\pi$'s end-node $v_n \in B$ and its end-node equal to $\pi$'s start-node $v_0 \in A$, so $\pi^* \in \mathcal{W}_{B \to A}$, and conversely by involution. (This resolves the wording-check item \texttt{walks\_same\_regardless\_of\_direction\_ambiguous} recorded in the workspace scratchpad for this row: walks are ordered sequences in the sense of def \ref{def:walks}, item~i., and the bijection between the two direction-sets is the explicit involution $\pi \leftrightarrow \pi^*$, not literal set equality.)

\medskip\noindent\textbf{(III) $\sigma$-blocking is invariant under reversal.}
We claim that, for every walk $\pi$ in $G$,
\[
    \pi \text{ is } C\text{-}\sigma\text{-blocked}
    \quad\iff\quad
    \pi^* \text{ is } C\text{-}\sigma\text{-blocked},
\tag{$\dagger$}
\]
both in the sense of def \ref{def:sigma_blocking}. By the involution $(\pi^*)^* = \pi$, the two directions of $(\dagger)$ are mutually each other's converses, so it suffices to prove the forward direction. We do so by tracing each of the two existential disjuncts of def \ref{def:sigma_blocking}, separately, through the position-shift $k \leftrightarrow n - k$ induced by reversal.

\smallskip\noindent\emph{Position-shift bookkeeping.}
By construction (II), the node at position $k$ on $\pi$ -- namely $v_k$ -- is the node at position $n - k$ on $\pi^*$ (which reads $\pi$'s node sequence backwards). The walk-incident edges at position $n - k$ on $\pi^*$ are slot $n - k - 1$ of $\pi^*$ (when $n - k \ge 1$) and slot $n - k$ of $\pi^*$ (when $n - k \le n - 1$). By the slot indexing of (II), slot $j$ of $\pi^*$ corresponds to slot $n - 1 - j$ of $\pi$; so:
\begin{itemize}
    \item slot $n - k - 1$ of $\pi^*$ corresponds to slot $n - 1 - (n - k - 1) = k$ of $\pi$ (i.e.\ to $a_k$, the right-incident slot at position $k$ on $\pi$), with the edge stored as $a_k^*$ per case (F$_E$)/(B$_E$)/(L) above;
    \item slot $n - k$ of $\pi^*$ corresponds to slot $n - 1 - (n - k) = k - 1$ of $\pi$ (i.e.\ to $a_{k-1}$, the left-incident slot at position $k$ on $\pi$), with the edge stored as $a_{k-1}^*$.
\end{itemize}
In words: the left-incident slot of $\pi^*$ at position $n - k$ is the reversal of the right-incident slot of $\pi$ at position $k$, and the right-incident slot of $\pi^*$ at position $n - k$ is the reversal of the left-incident slot of $\pi$ at position $k$. The walk-incident index sets at position $k$ on $\pi$ and at position $n - k$ on $\pi^*$ are therefore in canonical bijection (with the two slots swapped). Symbolically, abbreviating the walk-incident slots at position $k$ of $\pi$ as the typed pair $(s_{k-1}, s_k)$ (per the typed walk-step convention of \refrow{def_3_4}), the walk-incident slots at position $n - k$ of $\pi^*$ are $(s_k^*, s_{k-1}^*)$, where $s_j^*$ denotes the per-slot reversal of $s_j$.

\smallskip\noindent\emph{(III.a) Collider invariance.}
We show: position $k$ is a collider on $\pi$ (\refrow{def_3_15}, item~ii., under the side-aware reading committed to by the addition tag \texttt{[collider\_side\_aware\_walkstep\_predicates]}) iff position $n - k$ is a collider on $\pi^*$ in the same sense.

Per \refrow{def_3_15}, item~ii., position $k$ is a collider on $\pi$ iff $\mathrm{ah}_\pi(k) = 2$, i.e.\ both walk-incident edges at position $k$ -- $a_{k-1}$ (slot $k - 1$, present when $k \ge 1$) and $a_k$ (slot $k$, present when $k \le n - 1$) -- contribute an arrowhead at $v_k$. Under the side-aware reading the per-step contribution is read off the typed walk-step's constructor tag and walk-traversal direction (\refrow{def_3_15}, ``Encoding note: side-aware reading via a typed walk-step representation'' paragraph). The relevant predicates are the two per-step helpers:
\begin{itemize}
    \item ``$s$ contributes an arrowhead on its walk-traversal \emph{target} side'', written $\HeadAtTarget(s)$, which fires iff (a) $s$'s constructor tag is \texttt{forwardE} (a directed edge $(\text{source}, \text{target}) \in E$ traversed forwards has its head at the target), or (b) $s$'s constructor tag is \texttt{bidir} (a bidirected edge has heads at both endpoints), or (c) $s$'s constructor tag is \texttt{backwardE} \emph{and} the same stored pair simultaneously lies in $L$ (the writing-mirror re-firing disjunct $s(\text{source}, \text{target}) \in L$);
    \item ``$s$ contributes an arrowhead on its walk-traversal \emph{source} side'', written $\HeadAtSource(s)$, which is the symmetric predicate: it fires iff $s$'s constructor tag is \texttt{backwardE}, or \texttt{bidir}, or \texttt{forwardE} with the writing-mirror re-firing disjunct $s(\text{source}, \text{target}) \in L$.
\end{itemize}
At position $k$ on $\pi$ with walk-incident slots $(s_{k-1}, s_k)$, the node $v_k$ is the \emph{target} index of $s_{k-1}$ (the left-incident slot terminates at $v_k$) and the \emph{source} index of $s_k$ (the right-incident slot starts at $v_k$). Therefore the side-aware collider condition $\mathrm{ah}_\pi(k) = 2$ at an interior position $1 \le k \le n - 1$ reads, exactly,
\begin{equation}
    \IsCollider_\pi(k)
    \quad\iff\quad
    \HeadAtTarget(s_{k-1}) \;\land\; \HeadAtSource(s_k). \tag{$\ddagger$}
\end{equation}
(At end-positions $k \in \{0, n\}$ at most one of the two walk-incident slots exists, so $\mathrm{ah}_\pi(k) \le 1$ and the position is automatically a non-collider; the formula above is consistent with the side-aware classifier returning \emph{False} structurally at end-positions, since the relevant conjunct simply has no walk-incident slot to query.)

We need a per-step reversal-invariance lemma. Let $s : \WalkStep G u v$ be a typed walk-step (with walk-traversal source $u$ and target $v$), and let $s^* : \WalkStep G v u$ denote its per-slot reversal as constructed in (II). The constructor-tag flip from (II) gives, on each branch:
\begin{itemize}
    \item If $s$ is \texttt{forwardE} (stored pair $(u, v) \in E$), then $s^*$ is \texttt{backwardE} (stored pair $(u, v) \in E$, viewed at type indices $v \to u$ now). Then $\HeadAtTarget(s^*)$ (asking for a head at the new target $u$) reads, by the \texttt{backwardE} branch's clause (c)-mirror, the writing-mirror disjunct at the new step's source/target indices: $s(v, u) \in L$. By the canonical $\Sym_2$ identity $s(v, u) = s(u, v)$ (def \ref{def-cdmg}: $L$ is stored as unordered pairs), this equals $s(u, v) \in L$. On the original step's $\HeadAtSource(s)$ side, the \texttt{forwardE} branch's clause (c) reads $s(u, v) \in L$. \emph{Match}: $\HeadAtTarget(s^*) \iff \HeadAtSource(s)$. Symmetrically $\HeadAtSource(s^*)$ on the \texttt{backwardE} branch reads True; $\HeadAtTarget(s)$ on the \texttt{forwardE} branch reads True; match.
    \item If $s$ is \texttt{backwardE} (stored pair $(v, u) \in E$ in the original type-indexing $u \to v$, equivalently a directed edge $v \to u$ in $G$), then $s^*$ is \texttt{forwardE} (stored pair $(v, u) \in E$, viewed at type indices $v \to u$). Then $\HeadAtTarget(s^*)$ on the \texttt{forwardE} branch reads True; $\HeadAtSource(s)$ on the \texttt{backwardE} branch reads True; match. Symmetrically $\HeadAtSource(s^*)$ on the \texttt{forwardE} branch's clause (c) reads $s(v, u) \in L$, and $\HeadAtTarget(s)$ on the \texttt{backwardE} branch's clause (c)-mirror reads $s(u, v) \in L = s(v, u) \in L$ (by $\Sym_2$); match.
    \item If $s$ is \texttt{bidir} (stored witness $s(u, v) \in L$), then $s^*$ is \texttt{bidir} (stored witness $s(v, u) \in L = s(u, v) \in L$). All four readings ($\HeadAtTarget(s^*)$, $\HeadAtSource(s^*)$, $\HeadAtTarget(s)$, $\HeadAtSource(s)$) reduce to \emph{True} on their respective \texttt{bidir} branches; match.
\end{itemize}
Combining the three cases, we have the per-step reversal-invariance lemma:
\begin{equation}
    \HeadAtTarget(s^*) \;\iff\; \HeadAtSource(s)
    \qquad\text{and}\qquad
    \HeadAtSource(s^*) \;\iff\; \HeadAtTarget(s),
    \tag{rev}
\end{equation}
the equivalence in each constructor branch holding either by reflection ($\texttt{bidir}$, or the two ``always True'' clauses) or by the $\Sym_2$-canonical identity $s(u, v) = s(v, u)$ on the writing-mirror $L$-disjunct.

We now combine $(\ddagger)$, the position-shift bookkeeping, and (rev) to prove collider invariance. Fix a position $k$ on $\pi$.
\begin{itemize}
    \item \emph{Interior case $1 \le k \le n - 1$.} The walk-incident slots at position $k$ on $\pi$ are $(s_{k-1}, s_k)$; at position $n - k$ on $\pi^*$ they are $(s_k^*, s_{k-1}^*)$ (per the position-shift bookkeeping above: the left-incident slot of $\pi^*$ at position $n - k$ is the reversal of the right-incident slot $s_k$ of $\pi$ at position $k$, and the right-incident slot of $\pi^*$ at position $n - k$ is the reversal of the left-incident slot $s_{k-1}$ of $\pi$ at position $k$). The node $v_k$ on $\pi^*$ is the target index of the new left-incident slot $s_k^*$ and the source index of the new right-incident slot $s_{k-1}^*$ (just as it was target of $s_{k-1}$ and source of $s_k$ on $\pi$; the role-swap follows from the slot-swap under reversal). Therefore, by $(\ddagger)$ applied at position $n - k$ on $\pi^*$,
    \[
        \IsCollider_{\pi^*}(n - k)
        \;\iff\; \HeadAtTarget(s_k^*) \;\land\; \HeadAtSource(s_{k-1}^*)
        \;\overset{\mathrm{(rev)}}{\iff}\; \HeadAtSource(s_k) \;\land\; \HeadAtTarget(s_{k-1})
        \;\overset{\mathrm{And.comm}}{\iff}\; \HeadAtTarget(s_{k-1}) \;\land\; \HeadAtSource(s_k)
        \;\overset{(\ddagger)}{\iff}\; \IsCollider_\pi(k).
    \]
    \item \emph{End-position case $k \in \{0, n\}$.} Both ends map to each other under the position-shift ($k = 0$ on $\pi$ corresponds to $n - 0 = n$ on $\pi^*$ and vice versa). Both are non-collider end-positions on each respective walk by the LN's count bound $\mathrm{ah}_{\cdot}(k) \le 1$ at an end-position (\refrow{def_3_15}'s ``Arrowhead count at a position'' paragraph), so $\IsCollider_\pi(k)$ and $\IsCollider_{\pi^*}(n - k)$ are both False; the biconditional holds trivially.
\end{itemize}
Combining the two cases: position $k$ is a collider on $\pi$ iff position $n - k$ is a collider on $\pi^*$. By the side-aware ``non-collider iff arrowhead count $\le 1$'' classification (\refrow{def_3_15}, item~i.), which is the definitional negation of the side-aware collider predicate on the in-range fragment $\{0, 1, \dots, n\}$, the corresponding non-collider statement also follows:
\[
    \IsNonCollider_\pi(k) \;\iff\; \IsNonCollider_{\pi^*}(n - k).
\]

\smallskip\noindent\emph{Self-loop addendum (manager-accepted deviation \texttt{collider\_side\_aware\_at\_self\_loops}).}
At a walk-incident directed self-loop step $a_i = (v, v) \in E$ with $v_i = v_{i+1} = v$, the literal stored-pair test ``$a_i$ is an edge into $v_k$'' of def \ref{def-edge-relations}, item~ii.\ fires unconditionally at both adjacent positions $k \in \{i, i + 1\}$ for which $v_k = v$ (its $E$-clause's $e_2 = v_k$ condition reduces to $v = v$). The side-aware reading is strictly finer: at a self-loop step encoded as \texttt{forwardE} (type \texttt{WalkStep}$\,G\,v\,v$), $\HeadAtTarget$ on the \texttt{forwardE} branch reads True (the head is at the walk-traversal target side), while $\HeadAtSource$ reads $s(v, v) \in L$ via the writing-mirror disjunct. By \refrow{def_3_1}'s irreflexivity restriction \texttt{hL\_irrefl} on the bidirected channel ($\texttt{leanification/Chapter3\_GraphTheory/Section3\_1/CDMG.lean:376}$: $\forall \llbracket s : \Sym_2\,\Node \rrbracket,\ s \in L \to \neg\, s.\IsDiag$), the $\Sym_2$-diagonal witness $s(v, v)$ cannot lie in $L$. Hence the writing-mirror disjunct is \emph{vacuously false} at a self-loop, $\HeadAtSource$ on a \texttt{forwardE} self-loop step reduces to False, and the self-loop contributes exactly $1$ to $\mathrm{ah}_\pi(k)$ at its walk-traversal target position ($k = i + 1$) and $0$ at its walk-traversal source position ($k = i$). The same reading applies symmetrically to a self-loop encoded as \texttt{backwardE}, with the head moved to the source side.

Under reversal, the self-loop's constructor tag flips per (II): a \texttt{forwardE} self-loop step in $\pi$ at slot $i$ becomes a \texttt{backwardE} step in $\pi^*$ at slot $n - 1 - i$, and the $\HeadAtTarget$ / $\HeadAtSource$ counters at the two adjacent positions transpose accordingly. Both adjacent positions are the same physical node $v$ (since $v_i = v_{i+1} = v$), so the side-aware count $\mathrm{ah}_{\pi}(k) = \mathrm{ah}_{\pi^*}(n - k)$ remains balanced across reversal: whichever adjacent position carried the self-loop's $1$ on $\pi$ now carries the self-loop's $1$ on $\pi^*$ at the position-shifted index, with no double-counting and no artifact. In particular, the equality $\mathrm{ah}_\pi(k) = \mathrm{ah}_{\pi^*}(n - k)$ derived in the interior-case argument above goes through unchanged at self-loop positions: the per-step reversal-invariance lemma (rev) absorbs the self-loop case via the same constructor-tag-flip plus $\Sym_2$-canonical identity, with the writing-mirror disjunct vacuously false on both sides of the biconditional by \texttt{hL\_irrefl}. This is consistent with the canonical count of $1$ at a self-loop step recorded in \refrow{def_3_15}'s ``Treatment of directed self-loops'' paragraph.

\smallskip\noindent\emph{(III.b) Blockable non-collider invariance.}
We show: position $k$ is a blockable non-collider on $\pi$ (\refrow{def_3_16}, paragraph ``Blockable non-collider on $\pi$'', under the side-aware reading committed to by \refrow{def_3_15}'s addition tag) iff position $n - k$ is a blockable non-collider on $\pi^*$.

Per \refrow{def_3_16}, position $k$ is a blockable non-collider on $\pi$ iff (i)~it is a non-collider on $\pi$ (\refrow{def_3_15}, item~i., side-aware reading), AND (ii)~the blockable disjunction holds at $k$ on $\pi$:
\[
    k \in \{0, n\}
    \;\;\lor\;\;
    \HasBlockingLeftSlot_\pi(k)
    \;\;\lor\;\;
    \HasBlockingRightSlot_\pi(k),
\]
where the two helpers $\HasBlockingLeftSlot$ and $\HasBlockingRightSlot$ pattern-match on the walk-incident slot's constructor tag directly -- \emph{they do not consult any ``edge into $v_k$'' stored-pair test} -- and read $\HasBlockingLeftSlot_\pi(k) \iff a_{k-1} = (v_k, v_{k-1}) \in E \land v_{k-1} \notin \Sc^G(v_k)$ (slot $k - 1$ is a \texttt{backwardE} step with source $v_k$, target $v_{k-1}$, and the target lies outside $\Sc^G(v_k)$), respectively $\HasBlockingRightSlot_\pi(k) \iff a_k = (v_k, v_{k+1}) \in E \land v_{k+1} \notin \Sc^G(v_k)$ (slot $k$ is a \texttt{forwardE} step with source $v_k$, target $v_{k+1}$, and the target lies outside $\Sc^G(v_k)$). These two helpers are \emph{unchanged} between the original \refrow{def_3_16} encoding and the side-aware refactor (the refactor changes only the non-collider classifier referenced by clause (i)); they are net-zero under the refactor.

We show invariance of clauses (i) and (ii) separately.
\begin{itemize}
    \item \emph{Clause (i) -- non-collider status} -- is invariant under the position-shift $k \leftrightarrow n - k$ by (III.a): $\IsNonCollider_\pi(k) \iff \IsNonCollider_{\pi^*}(n - k)$, where both classifiers are the side-aware \refrow{def_3_15}, item~i.\ predicate. (At a position $k$ adjacent to a directed self-loop, the side-aware non-collider classifier returns True at both the position on $\pi$ and the position-shifted index on $\pi^*$, since the side-aware self-loop classification is the strict refinement noted in the self-loop addendum to (III.a) and the count is reversal-invariant; the same is true on walks with no self-loops, where the side-aware and stored-pair readings coincide.)
    \item \emph{Clause (ii) -- the blockable disjunction} -- we check each of the three disjuncts. Throughout, recall from the position-shift bookkeeping that the right-incident slot of $\pi$ at position $k$ (slot $k$, recorded edge $a_k$, constructor tag at that slot) corresponds under reversal to the left-incident slot of $\pi^*$ at position $n - k$ (slot $n - k - 1$, recorded edge $a_k^*$), with the constructor tag flipped per (II)'s rules (\texttt{forwardE} $\leftrightarrow$ \texttt{backwardE}; \texttt{bidir} $\leftrightarrow$ \texttt{bidir}). Symmetrically the left-incident slot of $\pi$ at $k$ becomes the right-incident slot of $\pi^*$ at $n - k$.
        \begin{itemize}
            \item \emph{End-position $k \in \{0, n\}$.} Under the position-shift $k \mapsto n - k$, the set $\{0, n\}$ is mapped to itself (with $0 \leftrightarrow n$), so $k \in \{0, n\}$ on $\pi$ iff $n - k \in \{0, n\}$ on $\pi^*$.
            \item \emph{Left-slot disjunct on $\pi$: $\HasBlockingLeftSlot_\pi(k)$.} Suppose $a_{k-1} = (v_k, v_{k-1}) \in E$ and $v_{k-1} \notin \Sc^G(v_k)$, i.e.\ the slot $k - 1$ of $\pi$ has constructor tag \texttt{backwardE} with stored pair $(v_k, v_{k-1}) \in E$, source $v_{k-1}$ and target $v_k$ at the slot's typed indices. Per (II) case (B$_E$), slot $k - 1$ of $\pi$ reverses to slot $n - 1 - (k - 1) = n - k$ of $\pi^*$ with constructor tag \texttt{forwardE}, stored pair $(v_k, v_{k-1}) \in E$, source $v_k$ and target $v_{k-1}$ at the reversed slot's typed indices. At position $n - k$ on $\pi^*$, slot $n - k$ is the right-incident slot; its constructor tag is \texttt{forwardE} and its source is $v_k$, so the right-slot helper at $\pi^*$, position $n - k$, fires precisely when its target $v_{k-1}$ lies outside $\Sc^G(v_k)$ -- which is the same condition as on $\pi$ (since $\Sc^G(\cdot)$ is a property of $G$ alone). Therefore $\HasBlockingLeftSlot_\pi(k) \iff \HasBlockingRightSlot_{\pi^*}(n - k)$.
            \item \emph{Right-slot disjunct on $\pi$: $\HasBlockingRightSlot_\pi(k)$.} Symmetric to the previous bullet via (II) case (F$_E$). Suppose $a_k = (v_k, v_{k+1}) \in E$ and $v_{k+1} \notin \Sc^G(v_k)$, i.e.\ slot $k$ of $\pi$ has constructor tag \texttt{forwardE} with stored pair $(v_k, v_{k+1}) \in E$, source $v_k$ and target $v_{k+1}$. Per (II) case (F$_E$), slot $k$ of $\pi$ reverses to slot $n - 1 - k$ of $\pi^*$ with constructor tag \texttt{backwardE}, stored pair $(v_k, v_{k+1}) \in E$, source $v_{k+1}$ and target $v_k$ at the reversed slot's typed indices. At position $n - k$ on $\pi^*$, slot $n - 1 - k = n - k - 1$ is the left-incident slot; its constructor tag is \texttt{backwardE} and its target (in the typed-walk-step indexing) is $v_k$, so the left-slot helper at $\pi^*$, position $n - k$, fires precisely when its source $v_{k+1}$ lies outside $\Sc^G(v_k)$ -- again the same condition as on $\pi$. Therefore $\HasBlockingRightSlot_\pi(k) \iff \HasBlockingLeftSlot_{\pi^*}(n - k)$.
        \end{itemize}
\end{itemize}
Combining the three disjuncts: the blockable disjunction at $k$ on $\pi$ (end-position OR left-slot OR right-slot) is equivalent to the blockable disjunction at $n - k$ on $\pi^*$ (end-position OR \emph{right}-slot OR \emph{left}-slot), which is the same three-disjunct disjunction with the two interior disjuncts swapped in order. The ``other endpoint $\notin \Sc^G(v_k)$'' clause attaches to the underlying graph-edge's other endpoint and to the node $v_k$ at the position, both of which transport verbatim under reversal; the membership predicate $\Sc^G(\cdot) = \Anc^G(\cdot) \cap \Desc^G(\cdot)$ (def \ref{def_3_5}, item~vii.) is a property of $G$ alone -- not of the walk -- and is therefore unaffected by reversal.

Combining: clause (i) and clause (ii) at $k$ on $\pi$ jointly transport, by (III.a) and the case-by-case check above, to clause (i) and clause (ii) at $n - k$ on $\pi^*$. Hence position $k$ is a (side-aware) blockable non-collider on $\pi$ iff position $n - k$ is a (side-aware) blockable non-collider on $\pi^*$.

\smallskip\noindent\emph{(III.c) $C$-$\sigma$-blocking invariance.}
We now show $(\dagger)$. Per def \ref{def:sigma_blocking}, $\pi$ is $C$-$\sigma$-blocked iff at least one of the following two existentials holds:
\begin{itemize}
    \item[(B1)] there exists a position $k \in \{0, 1, \dots, n\}$ such that $k$ is a collider on $\pi$ and $v_k \notin \Anc^G(C)$;
    \item[(B2)] there exists a position $k \in \{0, 1, \dots, n\}$ such that $k$ is a blockable non-collider on $\pi$ and $v_k \in C$.
\end{itemize}
Both ``collider on $\pi$'' and ``blockable non-collider on $\pi$'' are read in the side-aware sense of \refrow{def_3_15} / \refrow{def_3_16} above; the $\sigma$-blocked predicate of def \ref{def:sigma_blocking} consumes whichever reading the side-aware refactor commits to, and the two existentials above are the verbatim side-aware unfolding. We must show the same disjunction holds for $\pi^*$.

Suppose (B1) holds for $\pi$, with witness position $k$ and node $v_k$. By (III.a), position $n - k$ is a collider on $\pi^*$ (under the side-aware reading), and the node at position $n - k$ on $\pi^*$ is the same node $v_k$ (the node identity transports under reversal: $\pi^*$'s node-sequence is just $\pi$'s node-sequence read backwards). The condition $v_k \notin \Anc^G(C)$ is a condition on the node $v_k$ and on $G$ alone (def \ref{def_3_5}, item~iv.; the ancestor set in $G$ is independent of any walk), so it is unaffected by which walk we are considering. Hence position $n - k$ on $\pi^*$ is a collider with $v_k \notin \Anc^G(C)$, so (B1) holds for $\pi^*$ at position $n - k$.

Suppose (B2) holds for $\pi$, with witness position $k$ and node $v_k$. By (III.b), position $n - k$ is a blockable non-collider on $\pi^*$, and the node at position $n - k$ on $\pi^*$ is $v_k$. The condition $v_k \in C$ is a condition on the node $v_k$ alone, so it transports verbatim. Hence (B2) holds for $\pi^*$ at position $n - k$.

In either case, $\pi^*$ is $C$-$\sigma$-blocked. So $\pi$ $C$-$\sigma$-blocked $\Rightarrow$ $\pi^*$ $C$-$\sigma$-blocked; by the involution argument noted above, the converse follows, and $(\dagger)$ is established.

\medskip\noindent\textbf{(IV) The forward implication $(\star)$.}
Assume $A \sPerp_G B \given C$. By (I), this says: for every walk $\pi$ in $G$ with $v_0 \in A$ and $v_n \in B$, $\pi$ is $C$-$\sigma$-blocked.

Let $\pi' = (v_0', a_0', \dots, a_{n-1}', v_n')$ be an arbitrary walk in $G$ with $v_0' \in B$ and $v_n' \in A$. We must show $\pi'$ is $C$-$\sigma$-blocked. Apply the reversal construction of (II) to $\pi'$: by the bijection $\mathcal{W}_{B \to A} \to \mathcal{W}_{A \to B}$ (specialised at source-set $A$ and target-set $B$), $(\pi')^*$ is a walk in $G$ with start-node $v_n' \in A$ and end-node $v_0' \in B$, so $(\pi')^* \in \mathcal{W}_{A \to B}$. By the hypothesis $A \sPerp_G B \given C$ applied to $(\pi')^*$, the walk $(\pi')^*$ is $C$-$\sigma$-blocked. By (III) applied to $\pi'$ -- specifically, $(\dagger)$ applied to $\pi'$ gives $\pi'$ $C$-$\sigma$-blocked iff $(\pi')^*$ $C$-$\sigma$-blocked -- the walk $\pi'$ is $C$-$\sigma$-blocked. Since $\pi'$ was arbitrary, $B \sPerp_G A \given C$ holds, by the unfolding in (I).

\medskip\noindent\textbf{(V) Conclusion.}
By the reduction in (I), the biconditional $A \sPerp_G B \given C \iff B \sPerp_G A \given C$ follows from the forward implication $(\star)$, which is (IV); the converse implication is obtained by applying $(\star)$ to the pair $(B, A)$ in place of $(A, B)$, using the very same construction of (II) and invariance of (III), neither of which depends on the labelling of the two source/target sets. Hence
\[
    A \sPerp_G B \given C \iff B \sPerp_G A \given C,
\]
as required.
\end{proof}

\end{document}
